\section{Experiment, identification, and risk}

Throughout, \(n\ge1\) is the sample size, \(d\ge1\) is the number of baseline categories, and \(q\in(0,1]\) is a supplied lower bound on occupied-cell outcome arrival. Generic positive constants \(c\) and \(C\) are independent of these parameters unless dependence is stated explicitly. We write \(u\lesssim v\) for \(u\le Cv\) and \(u\asymp v\) for two such inequalities, uniformly over the stated parameter domain. Scalar errors are measured by absolute value. The arm index \leanref{sym:a}{\ensuremath{a}} and surrogate index \leanref{sym:s}{\ensuremath{s}} range over \(\{0,1\}\), while the baseline index \leanref{sym:x}{\ensuremath{x}} ranges over \(\{0,\ldots,d-1\}\). We use \leanref{sym:N}{\ensuremath{N}}, the effective sample-size scale, for \(N=nq\); it combines the sample size with the guaranteed within-cell arrival probability.

\subsection{Records and cells}

The experiment has a finite baseline category \leanref{sym:X}{\ensuremath{X}}, a binary treatment assignment \leanref{sym:A}{\ensuremath{A}}, binary potential surrogates \leanref{sym:S(a)}{\ensuremath{S(a)}}, binary potential outcomes \leanref{sym:Y(a)}{\ensuremath{Y(a)}}, and an outcome-arrival indicator \leanref{sym:R}{\ensuremath{R}}. The following definitions specify the record space and the coordinates selected by treatment assignment.

\begin{assumptionv}{ass:iid-sampling}[Independent and identically distributed sampling]
The observed records \(O_1,\ldots,O_n\) are independent and identically distributed according to the observed-data law induced by \(P\), the full-data probability law.
\end{assumptionv}

Treatment assignment selects the realized surrogate \leanref{sym:S}{\ensuremath{S}} from the two potential surrogate coordinates.

\begin{definitionv}{synth_2}[Full-data record coordinates]
A full-data record, for alphabet size \(d\in\mathbb N\), consists of a baseline category \(X\), a binary arm indicator \(A\), control and treated surrogate coordinates \(S(0)\) and \(S(1)\), control and treated outcome coordinates \(Y(0)\) and \(Y(1)\), and an arrival flag \(R\), with domains
\[
X\in\{x\in\mathbb Z:0\le x<d\},
\qquad
\bigl(A,S(0),S(1),Y(0),Y(1),R\bigr)\in\{0,1\}^{6}.
\]
\end{definitionv}

The assigned-arm outcome \leanref{sym:Y}{\ensuremath{Y}} is selected in the same way, and the finite record space carries its discrete measurable structure.

\begin{definitionv}{synth_3}[Potential and assigned surrogates]
The potential surrogate \(S(a)\), for \(a\in\{0,1\}\), denotes the record's binary control surrogate when \(a=0\) and binary treated surrogate when \(a=1\). The assigned-arm surrogate \(S\), for any full-data record with binary assignment \(A\in\{0,1\}\), is defined by
\[
S=
\begin{cases}
S(1), & A=1,\\
S(0), & A=0.
\end{cases}
\]
\end{definitionv}

A full-data probability law \leanref{sym:P}{\ensuremath{P}}, defined in \cref{obj:synth_1}, specifies the joint distribution of these coordinates.

\begin{definitionv}{synth_4}[Potential and assigned-arm outcomes]
The binary potential outcomes \(Y(0)\) and \(Y(1)\) are the control and treated outcome coordinates of a full-data record. For a nonnegative integer \(d\), such a record consists of a baseline category \(X\) in an alphabet of size \(d\) and binary coordinates \(A,S(0),S(1),Y(0),Y(1),R\). The assigned-arm outcome \(Y\) is defined by
\[
Y=
\begin{cases}
Y(1), & A=1,\\
Y(0), & A=0.
\end{cases}
\]
The space of full-data records is equipped with the measurable structure in which every subset is measurable.
\end{definitionv}

Outcome ascertainment is described within treatment--baseline--surrogate cells. The cell-index set \leanref{sym:\mathcal J_d}{\ensuremath{\mathcal J_d}} contains a cell \leanref{sym:j}{\ensuremath{j}} for each possible triple of these observed coordinates.

\begin{assumptionv}{ass:randomized-independence}[Independence of treatment assignment]
Under \(P\), the full-data probability law, treatment assignment \(A\) satisfies
\[
A\perp\!\!\!\perp\{X,S(0),S(1),Y(0),Y(1)\}.
\]
Here \(X\) is the baseline covariate, \(S(0),S(1)\) are the potential surrogates, and \(Y(0),Y(1)\) are the potential outcomes.
\end{assumptionv}

Two probabilities distinguish membership from outcome arrival: \leanref{sym:p_{axs}(P)}{\ensuremath{p_{axs}(P)}} is the full-membership cell probability, and \leanref{sym:\rho_{axs}(P)}{\ensuremath{\rho_{axs}(P)}} is the conditional arrival probability in an occupied cell.

\begin{assumptionv}{ass:balanced-randomization}[Balanced treatment assignment]
Under \(P\), the full-data probability law, treatment assignment \(A\) satisfies
\[
P(A=1)=\frac12.
\]
\end{assumptionv}

Thus an occupied cell has positive membership probability, and its arrived-cell probability is the product of its membership and conditional arrival probabilities. For the treatment contrast, the treatment sign \leanref{sym:g(a)}{\ensuremath{g(a)}} assigns opposite signs to the two arms.

\begin{definitionv}{synth_1}[Full-data probability law]
The full-data law \(P\), for a nonnegative integer \(d\), is a probability measure on the finite record space
\[
(X,A,S(0),S(1),Y(0),Y(1),R)
\in \{0,\ldots,d-1\}\times\{0,1\}^{6}.
\]
Here \(X\) is the baseline category, \(A\) is the treatment arm, \(S(0)\) and \(S(1)\) are the potential surrogates, \(Y(0)\) and \(Y(1)\) are the potential outcomes, and \(R\) is the arrival flag. The baseline alphabet is empty when \(d=0\).
\end{definitionv}

The observed record \leanref{sym:O}{\ensuremath{O}} is \((X,A,S,R,RY)\). Its sampled copies \leanref{sym:O_i}{\ensuremath{O_i}} retain cell membership and arrival status for every individual; when \(R=1\), the final coordinate records the outcome. The canonical observed sample law \leanref{sym:\mathbf O_n}{\ensuremath{\mathbf O_n}}, defined in \cref{obj:synth_5}, is the corresponding product distribution.

\begin{assumptionv}{ass:surrogate-consistency}[Consistency of the surrogate]
The realized surrogate \(S\) equals the potential surrogate corresponding to treatment assignment \(A\):
\[
S=S(A)\qquad\text{almost surely under }P,
\]
where \(P\) is the full-data probability law.
\end{assumptionv}

\subsection{Sampling, randomization, and outcome arrival}

Outcome consistency completes the link between observed and potential outcomes.

\begin{assumptionv}{ass:outcome-consistency}[Consistency of the outcome]
The realized outcome \(Y\) equals the potential outcome corresponding to treatment assignment \(A\):
\[
Y=Y(A)\qquad\text{almost surely under }P,
\]
where \(P\) is the full-data probability law.
\end{assumptionv}

Independent and identically distributed sampling is imposed within each triangular-array row, allowing the underlying law and alphabet size to vary across experiments \citep{ZengBalakrishnanHanKennedy2026}. Outcome arrival is governed by the following conditional independence condition.

\begin{assumptionv}{ass:arrival-mar}[Conditionally independent outcome arrival]
Under \(P\), the full-data probability law, the outcome-arrival indicator \(R\) and realized outcome \(Y\) are conditionally independent given baseline \(X\), treatment assignment \(A\), and realized surrogate \(S\):
\[
R\perp\!\!\!\perp Y\mid(X,A,S).
\]
\end{assumptionv}

For later aggregation, we index treatment--baseline--surrogate cells explicitly.

\begin{definitionv}{synth_6}[Cell alphabet and membership]
The cell alphabet \(\mathcal J_d\), for a nonnegative integer \(d\), is
\[
\mathcal J_d
=
\{0,1\}\times\{x\in\mathbb Z:0\le x<d\}\times\{0,1\}.
\]
For a full-data record with baseline category \(X\), binary assignment \(A\), and realized surrogate \(S\), membership in a cell \(j=(a,x,s)\in\mathcal J_d\) is defined by
\[
A=a \quad\land\quad X=x \quad\land\quad S=s.
\]
\end{definitionv}

The next definition records each cell's population mass and its conditional arrival probability.

\begin{definitionv}{synth_16}[Cell probabilities \(p_{axs}(P)\) and \(\rho_{axs}(P)\)]
The full-membership probability of a cell is
\[
p_{axs}(P):=P(A=a,X=x,S=s),
\qquad
(a,x,s)\in\mathcal J_d
=\{0,1\}\times\{0,\ldots,d-1\}\times\{0,1\},
\]
where \(P\) is a full-data law. On every occupied cell, meaning \(p_{axs}(P)>0\), define its conditional outcome-arrival probability by
\[
\rho_{axs}(P)
:=
P(R=1\mid A=a,X=x,S=s)
=
\frac{P(A=a,X=x,S=s,R=1)}{p_{axs}(P)}.
\]
\end{definitionv}

The supplied arrival floor is imposed only on occupied cells.

\begin{assumptionv}{ass:occupied-cell-arrival}[Occupied-cell arrival floor]
For each cell
\[
(a,x,s)\in\mathcal J_d
=\{0,1\}\times\{0,\ldots,d-1\}\times\{0,1\},
\]
let \(p_{axs}(P)\), the full-membership cell probability, and \(\rho_{axs}(P)\), the conditional outcome-arrival probability on an occupied cell, be defined as in \cref{obj:synth_16}. The occupied-cell arrival condition is
\[
p_{axs}(P)>0
\quad\Longrightarrow\quad
\rho_{axs}(P)\ge q.
\]
\end{assumptionv}

The preceding randomization, consistency, and arrival restrictions define the full positive-arrival model.

\begin{definitionv}{def:unrestricted-arrival-model-class}[Arrival-floor model \(\mathcal M^{+}_{n,d,q}\)]
The class \(\mathcal M^{+}_{n,d,q}\) consists of full-data laws \(P\), for integers \(n,d\ge1\) and an arrival floor \(0<q\le1\), satisfying the following conditions:
\begin{itemize}
\item \textbf{(Randomized independence.)} \cref{obj:ass:randomized-independence}.
\item \textbf{(Balanced randomization.)} \cref{obj:ass:balanced-randomization}.
\item \textbf{(Surrogate consistency.)} \cref{obj:ass:surrogate-consistency}.
\item \textbf{(Outcome consistency.)} \cref{obj:ass:outcome-consistency}.
\item \textbf{(Arrival MAR.)} \cref{obj:ass:arrival-mar}.
\item \textbf{(Occupied-cell arrival.)} \cref{obj:ass:occupied-cell-arrival}.
\end{itemize}
The arrival floor ranges over the entire interval \(0<q\le1\). Independent, identically distributed sampling is imposed separately on the experiment through \cref{obj:ass:iid-sampling}.
\end{definitionv}

Outcome arrival missing at random given treatment, baseline, and surrogate equates the outcome distribution among arrivals with the full-cell outcome distribution whenever arrivals have positive probability \citep{Rubin1976,Robins1994,KallusMao2025}. The surrogate therefore serves as an observed conditioning variable for ascertainment: its realized value can capture variation in arrival associated with the assigned treatment and the outcome. Positive ascertainment within occupied cells is imposed next.

\begin{definitionv}{def:ate-functional}[Average treatment effect \(\tau(P)\)]
The average treatment effect is
\[
\tau(P):=E_P[Y(1)-Y(0)],
\]
where \(Y(1)\) and \(Y(0)\) are the potential outcomes under treatment and control.
\end{definitionv}

The occupied-cell arrival floor is specific to this analysis. It guarantees an outcome-arrival probability of at least the supplied value \(q\) in every cell represented in the population, while permitting both cell membership probabilities and arrival probabilities to vary across cells. This condition supplies positive arrived-cell mass wherever the population cell contribution can be nonzero. A prospective analysis must justify this lower bound; the stated tuning and guarantees apply when it is valid.

Together, the randomization, consistency, and arrival restrictions define the arrival-floor model \leanref{sym:\mathcal M^{+}_{n,d,q}}{\ensuremath{\mathcal M^{+}_{n,d,q}}} in \cref{obj:def:unrestricted-arrival-model-class}.

\begin{definitionv}{synth_19}[Treatment sign function]
Define the treatment sign function \(g:\{0,1\}\to\{-1,1\}\) by
\[
g(a):=2a-1,\qquad a\in\{0,1\}.
\]
Thus \(g(0)=-1\) and \(g(1)=1\).
\end{definitionv}

This family covers the full range of positive arrival floors, including complete outcome arrival at \(q=1\). Sampling determines the product experiment induced by each member of the family.

\subsection{Identification and squared-error risk}

The causal target is the population average treatment effect \leanref{sym:\tau(P)}{\ensuremath{\tau(P)}} in \cref{obj:def:ate-functional}.

\begin{definitionv}{def:cell-functional}[Signed cell functional \(\Phi(P)\)]
The total signed observed-cell functional is
\[
\Phi(P):=2\sum_{(a,x,s)\in\mathcal J_d}g(a)c_{axs}(P),
\]
where \(\mathcal J_d\) is the treatment–baseline–surrogate cell-index set and \(g(a)\) is the treatment sign function. On this index set, with \(x\in\{0,\ldots,d-1\}\), the full-membership cell probability \(p_{axs}(P)\) is defined in \cref{obj:synth_16}.

For a cell with positive arrived-cell probability,
\[
P(A=a,X=x,S=s,R=1)>0,
\]
define its arrived-cell outcome mean and weighted contribution by
\[
\mu_{axs}(P):=E_P[Y\mid A=a,X=x,S=s,R=1],
\qquad
c_{axs}(P):=p_{axs}(P)\mu_{axs}(P).
\]
When the arrived-cell probability is zero, including whenever \(p_{axs}(P)=0\), set
\[
\mu_{axs}(P):=0,
\qquad
c_{axs}(P):=0.
\]
Thus \(\Phi\) is defined on every full-data law.
\end{definitionv}

Binary potential outcomes place this target in \([-1,1]\). Its observed-data representation combines full-membership probabilities with outcome means among arrivals. The total signed cell functional \leanref{sym:\Phi(P)}{\ensuremath{\Phi(P)}}, defined in \cref{obj:def:cell-functional}, uses an arrived-cell outcome mean \leanref{sym:\mu_{axs}(P)}{\ensuremath{\mu_{axs}(P)}} and a membership-weighted contribution \leanref{sym:c_{axs}(P)}{\ensuremath{c_{axs}(P)}}, with explicit conventions for cells having zero arrived mass.

\begin{definitionv}{synth_5}[Observed sample and its canonical law]
The random observed sample is \(\mathbf O_n=(O_1,\ldots,O_n)\), where \(O_i=(X_i,A_i,S_i,R_i,R_iY_i)\). For nonnegative integers \(n,d\) and a full-data probability law \(P\) with \(d\) baseline categories, its canonical law is
\[
\mathcal L_P(\mathbf O_n):=\bigl(\mathcal L_P(X,A,S,R,RY)\bigr)^{\otimes n}.
\]
Here \(\mathcal L_P(X,A,S,R,RY)\) is the distribution under \(P\) of the observed record \((X,A,S,R,RY)\), whose last four coordinates are binary and whose final coordinate is the product of arrival and outcome. Thus \(\mathcal L_P(\mathbf O_n)\) is the law of \(n\) independent observed records with this common distribution.
\end{definitionv}

To distinguish the two roles in subsequent formulas, we write \(\mathbf O_n=(O_1,\ldots,O_n)\) for the random observed-record tuple and \(\mathcal L_P(\mathbf O_n)\) for its product law. The canonical-law display above specifies the latter distribution.

For every law in \cref{obj:def:unrestricted-arrival-model-class}, \cref{obj:lem:unrestricted-cell-identification} establishes \(\tau(P)=\Phi(P)\); its proof appears in the appendix. The representation assigns zero contribution to null cells and uses observable arrived-cell means on occupied cells. Full cell membership remains observed even when the outcome is missing, so the representation separates population cell weights from within-cell outcome ascertainment. Balanced assignment supplies the factor of two in the signed contrast.

We evaluate estimation over the entire arrival-floor model through the minimax squared-error risk \leanref{sym:\mathfrak R^{+}_{n,d,q}}{\ensuremath{\mathfrak R^{+}_{n,d,q}}} in \cref{obj:def:unrestricted-minimax-risk}. Its decision class \leanref{sym:\mathcal T_n}{\ensuremath{\mathcal T_n}} contains all parameter-independent bounded statistical decision rules, including a possibly randomized estimator \leanref{sym:\mathsf T}{\ensuremath{\mathsf T}}.

\begin{definitionv}{def:unrestricted-minimax-risk}[Unrestricted minimax squared-error risk]
The minimax squared-error risk \(\mathfrak R^{+}_{n,d,q}\), for \(n,d\in\mathbb N\) and \(q\in\mathbb R\), is defined by
\[
\mathfrak R^{+}_{n,d,q}
:=
\inf_{\mathsf T\in\mathcal T_n}
\sup_{P\in\mathcal M^{+}_{n,d,q}}
\int\!\int_{[-1,1]}
(h-\tau(P))^2\,
\mathsf T(o,dh)\,
\mathcal L_P(\mathbf O_n)(do).
\]
Here \(\mathcal M^{+}_{n,d,q}\) is the full-data model class in \cref{obj:def:unrestricted-arrival-model-class}, and \(\tau(P)\) is the population average treatment effect in \cref{obj:def:ate-functional}. The class \(\mathcal T_n\) consists of all parameter-independent Markov kernels from the observed sample space to \([-1,1]\), including deterministic estimators as Dirac kernels.

The random observed-record tuple \(\mathbf O_n\) and its product law under the sampling in \cref{obj:ass:iid-sampling} are
\[
\mathbf O_n=(O_1,\ldots,O_n),
\qquad
\mathcal L_P(\mathbf O_n)
=
\bigl(\mathcal L_P(O_1)\bigr)^{\otimes n},
\]
where each \(O_i\) is an observed record induced by \(P\). The risk integrates over both the observed sample and the kernel draw.

By \cref{obj:lem:randomized-kernel-barycenter}, the defining infimum equals the infimum over deterministic estimators taking values in \([-1,1]\). At \(q=1\), the same decision class \(\mathcal T_n\) defines the minimax risk over \(\mathcal M^{+}_{n,d,1}\), whose laws satisfy \(R=1\) almost surely.
\end{definitionv}

The inner integral averages over the procedure's conditional output distribution given the sample, and the outer integral averages over sampling. Parameter independence means that the procedure is common to all laws in the specified model. Under squared loss, replacing a randomized output by its conditional mean given the sample weakly reduces risk for each law, as stated in \cref{obj:lem:randomized-kernel-barycenter}. This comparison makes the all-procedure criterion equivalent to its deterministic counterpart.

\subsection{The rare-arrival experiment and comparison rate}

A second model family restricts the arrival floor relative to the effective sample-size scale. Its defining condition uses the logarithmic scale \leanref{sym:\ell}{\ensuremath{\ell}}, which equals \(\log(e+N)\) on the statistical parameter domain fixed above.

\begin{definitionv}{synth_7}[Total real logarithmic scale]
The logarithmic scale \(\ell\), for \(n\in\mathbb N\) and \(q\in\mathbb R\), is defined by
\[
\ell=
\begin{cases}
\log\bigl|\exp(1)+nq\bigr|, & \exp(1)+nq\ne 0,\\
0, & \exp(1)+nq=0,
\end{cases}
\]
where \(\log\) on positive arguments is the natural logarithm.
\end{definitionv}

The rare-arrival regime is selected by the following restriction on \(q\) and this logarithmic scale.

\begin{assumptionv}{ass:rare-arrival-slice}[Rare-arrival restriction]
The arrival floor \(q\) and logarithmic scale \(\ell\) satisfy
\[
q\ell^2\le \frac{1}{64}.
\]
\end{assumptionv}

The rare-arrival restriction is specific to this analysis and selects parameter triples for which the arrival floor is small relative to the squared logarithmic scale. It restricts the guaranteed floor; individual occupied cells may have larger arrival probabilities. The rare-arrival model \leanref{sym:\mathcal M_{n,d,q}}{\ensuremath{\mathcal M_{n,d,q}}} in \cref{obj:def:model-class} combines this parameter restriction with the preceding statistical conditions.

\begin{definitionv}{def:model-class}[Rare-arrival model \(\mathcal M_{n,d,q}\)]
\(\mathcal M_{n,d,q}\) is the class of full-data probability laws \(P\) satisfying the following conditions:
\begin{itemize}
\item \textbf{(Parameters.)} \(n,d\ge1\) and \(0<q\le1\).
\item \textbf{(Sampling.)} \cref{obj:ass:iid-sampling} holds.
\item \textbf{(Randomization.)} \cref{obj:ass:randomized-independence,obj:ass:balanced-randomization} hold.
\item \textbf{(Consistency.)} \cref{obj:ass:surrogate-consistency,obj:ass:outcome-consistency} hold.
\item \textbf{(Arrival.)} \cref{obj:ass:arrival-mar,obj:ass:occupied-cell-arrival} hold.
\item \textbf{(Rare-arrival restriction.)} \cref{obj:ass:rare-arrival-slice} holds.
\end{itemize}
\end{definitionv}

At parameter triples satisfying \cref{obj:ass:rare-arrival-slice}, this model uses the same full-data restrictions as \cref{obj:def:unrestricted-arrival-model-class}, together with the stated sampling condition. The observed experiment \leanref{sym:\mathcal E_{n,d,q}}{\ensuremath{\mathcal E_{n,d,q}}} in \cref{obj:def:observed-experiment} is explicitly indexed by this rare-arrival model.

\begin{definitionv}{def:observed-experiment}[Observed sample experiment]
The observed experiment \(\mathcal E_{n,d,q}\), for \(n,d\in\mathbb N\) and \(q\in\mathbb R\), is the family
\[
\mathcal E_{n,d,q}
:=
\bigl\{\mathcal L_P(\mathbf O_n):P\in\mathcal M_{n,d,q}\bigr\},
\]
where \(\mathcal M_{n,d,q}\) is the model class defined in \cref{obj:def:model-class}. Here \(\mathbf O_n\), the random observed-record tuple, and its law are specified by
\[
\mathbf O_n=(O_i)_{i=1}^{n},
\qquad
O_i=(X_i,A_i,S_i,R_i,R_iY_i),
\qquad
\mathcal L_P(\mathbf O_n)
=
\bigl(\mathcal L_P(X,A,S,R,RY)\bigr)^{\otimes n}.
\]
Thus the records are independent and identically distributed with the observed marginal law induced by \(P\). Each record retains the baseline category \(X\in\{0,\ldots,d-1\}\), assignment \(A\), realized surrogate \(S\), arrival indicator \(R\), and arrival-multiplied outcome \(RY\); the last four coordinates are binary. The record space carries the discrete sigma-algebra.
\end{definitionv}

The decision problem retains every sampled cell membership, including those whose outcomes have yet to arrive. Its minimax squared-error risk \leanref{sym:\mathfrak R_{n,d,q}}{\ensuremath{\mathfrak R_{n,d,q}}}, defined in \cref{obj:def:minimax-risk}, uses the same decision class and joint sample-and-procedure expectation as the unrestricted criterion.

\begin{definitionv}{def:minimax-risk}[Minimax squared-error risk]
The minimax squared-error risk \(\mathfrak R_{n,d,q}\), for \(n,d\in\mathbb N\) and \(q\in\mathbb R\), is
\[
\mathfrak R_{n,d,q}
:=
\inf_{\mathsf T\in\mathcal T_n}
\sup_{P\in\mathcal M_{n,d,q}}
\int\!\int_{[-1,1]}
(h-\tau(P))^2\,
\mathsf T(o,dh)\,
\mathcal L_P(\mathbf O_n)(do).
\]
Here \(\mathcal M_{n,d,q}\) is the full-data model class of \cref{obj:def:model-class}, and \(\tau(P)\) is the population average treatment effect of \cref{obj:def:ate-functional}. The observed sample \(\mathbf O_n=(O_i)_{i=1}^n\) denotes the random tuple of observed records. Its law, under \cref{obj:ass:iid-sampling}, is
\[
\mathcal L_P(\mathbf O_n)
=
\bigotimes_{i=1}^{n}\mathcal L_P(O_i).
\]
The decision class \(\mathcal T_n\) consists of all parameter-independent Markov kernels from the observed-sample space to \([-1,1]\), including deterministic estimators represented by Dirac kernels.

The notation \(E_P[(\mathsf T-\tau(P))^2]\) denotes the displayed joint sample-and-kernel risk. By \cref{obj:lem:randomized-kernel-barycenter}, the minimax value also equals the infimum of the maximal squared-error risk over deterministic estimators taking values in \([-1,1]\).
\end{definitionv}

The two risk definitions distinguish the parameter domains over which the subsequent comparisons are stated. For the rare-arrival analysis, the comparison-rate expression \leanref{sym:r_{n,d,q}}{\ensuremath{r_{n,d,q}}} in \cref{obj:def:frontier-rate} combines the effective sample-size scale with the number of baseline categories.

\begin{definitionv}{def:frontier-rate}[Frontier rate \(r_{n,d,q}\)]
The frontier rate is
\[
r_{n,d,q}
:=
\min\left\{
1,\,
\frac{1}{nq}
+
\left[\frac{d}{nq\log(e+nq)}\right]^2
\right\}.
\]
\end{definitionv}

In terms of the setup notation, the expression is \(\min\{1,N^{-1}+[d/(N\ell)]^2\}\). Its first term records the reciprocal effective sample-size scale, while its second term expresses the joint dependence on alphabet size and outcome ascertainment. The truncation at one matches the uniform squared-error bound of the estimator that always reports zero for a target in \([-1,1]\). The estimation results that follow compare the defined risks with these scales under their stated ascertainment regimes.
